\documentstyle[%


righttag%


,11pt%


,amssymb%


,russian%               remove in Eng.


,izvru%                 Set izven% in Eng.


]{amsart}





% 





\newcommand{\la}{\langle}


\newcommand{\ra}{\rangle}


\newcommand{\diag}{\operatorname{diag}}


\newcommand{\sign}{\operatorname{sign}}


\newcommand{\Int}{\operatorname{int}}


\newcommand{\tts}[1]{}





\theoremstyle{plain}


\theorembodyfont{\it}


\newtheorem{thms}{\hskip\parindent Теорема}[section]





\theoremstyle{definition}


\theorembodyfont{\rm}


\newtheorem{defns}{\hskip\parindent Определение}[section]


\numberwithin{equation}{section}





%\hoffset=-15mm%                  For Eng.


\setcounter{page}{75}


\god{2004}


\nomer{12~(511)} %                        Remove in Eng.


%\nomer{Vol.\,48, No.\,12}%               Set in Eng.


%\str{\pageref{begin}--\pageref{end}}%  Set in Eng.


\udk{517.956}





\author{\it В.А.\,Терлецкий}


% NB:   ^^ remove in Eng.


\title{Обобщенное решение одномерных полулинейных


гиперболических систем со смешанными условиями}


\thanks{Работа выполнена при финансовой поддержке Российского фонда


фундаментальных исследований (проект \No\,02-01-00243), 


Министерства образования Российской Федерации, 


грант \No\,Е02-1.0-60, и программы ``Университеты России''


(проект \No\,УР.03.01.002).}





\date{}


%\pagestyle{myheadings}%         Set in Eng.


\pagestyle{plain} %              Remove in Eng.


\begin{document}


%\thispagestyle{plain}%          Set in Eng.


\maketitle


\label{begin}








Системы гиперболических уравнений имеют многочисленные (см., напр., [1]--[4])


приложения. При этом в приложениях зачастую приходится 


сталкиваться с негладкими данными, входящими либо в правую часть системы,


либо в ее начально-граничные условия, что приводит к необходимости расширять


запас гладких функций, которому только лишь и могут принадлежать классические


решения.  Проблема конструирования обобщенного решения включает в себя не 


только обоснование его корректности, но и выявление аналитических свойств, 


а также построение оценок роста относительно входных данных. Изучение


этих вопросов с различных позиций проводилось, например, в [1], [5]--[7].





В предлагаемой работе обобщенное решение определяется как решение системы 


интегральных уравнений, эквивалентной смешанной дифференциальной задаче на 


гладких решениях. В отличие от [1], [5]--[7] 


гиперболическая система рассматривается 


в общей, неинвариантной форме, а аппарат метода характеристик и инвариантов 


Римана задействован в качестве основного инструмента исследований. Смешанные 


условия конструируются в нелинейном виде. Доказываются существование и 


единственность обобщенного решения $x$ в предположении, что входные параметры 


задачи удовлетворяют условию Липшица по $x$ и суммируемы по Лебегу с некоторой 


степенью $p \geq 1$ по независимым переменным. Строится равномерная, 


справедливая  почти всюду в области независимых переменных оценка приращения


решения относительно возмущений входных данных. Обосновывается абсолютная


непрерывность линейных комбинаций компонент обобщенного решения (инвариантов


Римана) вдоль характеристик соответствующих семейств и формула интегрирования


по частям.








\section{Постановка задачи}





Пусть в плоскости переменных $(s,t)$ границей прямоугольника 


${\Pi} =S \times T$, $S=(s_0, s_1)$, $T=(t_0,t_1)$, является


$\partial {\Pi}=\overset{\vee}{\partial}{}{\Pi} \cup 


\overset{\wedge}{\partial}{}{\Pi} \cup \overset{<}{\partial}{}{\Pi}


\cup \overset{>}{\partial}{\Pi}$, где 


$\overset{<}{\partial}{}{\Pi}=\{(s,t): s=s_0$, $t \in T \}$,


$\overset{>}{\partial}{}{\Pi}=\{(s,t): s=s_1$, $t \in T \}$,


$\overset{\vee}{\partial}{}{\Pi}=\{(s,t): s \in S$, $t=t_0 \}$,


$\overset{\wedge}{\partial}{}{\Pi}=\{(s,t): s \in S$, $t=t_1 \}$.


$\overset{<>}{\partial}{}{\Pi}=\overset{<}{\partial}{}{\Pi} \cup


\overset{>}{\partial}{}{\Pi}$ --- боковая граница.


В прямоугольнике $\Pi$ определим систему дифференциальных уравнений 


с частными производными 


\begin{equation}


x_t+A(s,t)x_s=f(x,s,t).


\end{equation}


Здесь $x=x(s,t)$, $x(s,t) \in R^n$, --- искомое решение, $A=A(s,t)$, 


$A(s,t) \in R^{n \times n}$ --- заданная матричная функция, а


$f=f(x,s,t)$, $f(x,s,t) \in R^n$ --- заданная вектор-функция.





\begin{defns}


Систему дифференциальных уравнений (1.1) и 


дифференциальный оператор 


\begin{equation}


Dx=x_t+A(s,t)x_s 


\end{equation}


будем называть гиперболическими, если все собственные значения 


$\lambda_i=\lambda_i(s,t)$, $i=1,2,\dotsc,n$, матрицы $A$ являются


вещественными в $\Pi$ функциями, причем их алгебраическая и 


геометрическая кратности совпадают. 


\end{defns}





Как известно [8], в этом случае существуют левые 


$\ell^{(i)}=\ell^{(i)}(s,t)$ и правые 


$p^{(i)}=p^{(i)}(s,t)$, $i=1,2,\dotsc,n$, собственные векторы матрицы $A$ 


такие, что матрицы


${\cal L}=(\ell^{(1)}, \ell^{(2)},\dotsc,\ell^{(n)})$, \linebreak


${\cal P}=(p^{(1)},p^{(2)},\dotsc,p^{(n)})$ удовлетворяют соотношениям


\begin{equation}


{\cal L}^{\prime}(s,t) A(s,t)=\Lambda(s,t) {\cal L}^{\prime}(s,t), 


\quad A(s,t) {\cal P}(s,t)={\cal P}(s,t) \Lambda(s,t), 


\quad {\cal L}^{\prime}(s,t) {\cal P}(s,t)=\mathrm E.


\end{equation}


Здесь $\Lambda=\diag \{\lambda_1, \lambda_2,\dotsc,\lambda_n\}$, $\prime$ ---


знак транспонирования, $\mathrm E$ --- единичная матрица.





Дополнительно предположим, что матричные функции $\Lambda$, ${\cal L}$ и 


$\cal P$ непрерывны и непрерывно дифференцируемы в прямоугольнике 


$\ov {\Pi}={\Pi} \cup \partial {\Pi}$, а собственные значения 


$\lambda_i$ матрицы $\mathrm A$ удовлетворяют условию: либо 


$\lambda_i(s,t)=\lambda_j(s,t)$, либо $\lambda_i(s,t) \neq \lambda_j(s,t)$


при произвольных $i, j=1,2,\dotsc,n$, $i \neq j$, и всех $(s,t) \in \ov{\Pi}$.





Будем считать вектор-функцию $f(x,s,t)$ непрерывной по Липшицу по переменной


$x \in R^n$ при фиксированных $(s,t) \in {\Pi}$ и интегрируемой по 


Лебегу в $\Pi$ при фиксированных $x \in R^n$.





При сделанных предположениях относительно дифференциального оператора (1.2)


и вектор-функции $f(x,s,t)$ требуется определить корректный способ постановки 


начально-краевых условий, построить обобщенное решение системы (1.1), 


обосновать его существование и единственность, а также установить


оценки роста и аналитические свойства этого решения.











\section{Характеристики гиперболического оператора и инварианты Римана}





В этом параграфе для системы (1.1) выпишем так называемую инвариантную 


систему, воспользовавшись широко известным методом характеристик


(напр., [1], [2]), берущим свое начало от работы [9].





Введем характеристики гиперболического оператора (1.2) 


с помощью обыкновенных дифференциальных уравнений


\begin{equation}


\frac{ds}{dt}=\lambda_i(s,t),\quad i=1,2,\dotsc,n.


\end{equation}


Пусть $s=s^{(i)}(\xi,\tau;t)$ --- решение $i$-го уравнения (2.1), 


проходящее через точку $(\xi,\tau) \in {\ov{\Pi}}$. В силу гладкости 


собственных значений $\lambda_i(s,t)$ такие функции существуют, единственны, 


непрерывны и непрерывно дифференцируемы в прямоугольнике $\Pi$. Таким 


образом, каждое уравнение (2.1) порождает в $\Pi$ семейство интегральных 


кривых $s=s^{(i)}(\xi,\tau;t)$, которые называются характеристиками $i$-го 


семейства характеристик системы (1.1) (оператора (1.2)).





Отметим простое, но полезное обстоятельство. Одну и ту же интегральную кривую


$s=s^{(i)}(\xi,\tau;t)$ можно одновременно рассматривать и как решение


$t=t^{(i)}(\xi,\tau;s)$ дифференциального уравнения


$$


\frac{dt}{ds}=\frac{1}{\lambda_i(s,t)}


$$


в тех точках $(s,t) \in {\Pi}$, где $\lambda_i(s,t) \neq 0$, $i=1,2,\dotsc,n$.





Определим теперь инварианты Римана $r=r(s,t)$, $r(s,t) \in R^n$,


воспользовавшись линейной невырожденной заменой переменных


\begin{equation}


r(s,t)={\cal L}^{\prime}(s,t) x(s,t), 


\quad x(s,t)={\cal P}(s,t) r(s,t).


\end{equation}





Предположим, что вектор-функция $x=x(s,t)$ является классическим, т.\,е.


непрерывным и гладким, решением системы (1.1), и построим соответствующую 


систему дифференциальных уравнений, которой должны удовлетворять инварианты


Римана. Для этого подставим в (1.1) решение $x$ в виде (2.2). Будем иметь


$$


f({\cal P}r,s,t)=({\cal P}r)_t+A({\cal P}r)_s=D{\cal P}\cdot r +


{\cal P}(r_t+{\cal L}^{\prime}A{\cal P}r_s).


$$


Отсюда и из (1.3) получим систему для инвариантов Римана в виде


\begin{equation}


r_t+\Lambda(s,t)r_s={\cal L}^{\prime} (f({\cal P}r,s,t) - 


D{\cal P}\cdot r).


\end{equation}





Известна и другая форма записи инвариантной системы, соответствующей системе


(1.1). К ней можно прийти обратным путем, вычислив значение оператора


$$


D_{\Lambda} r=r_t+\Lambda(s,t) r_s.


$$


Действительно,


$$


D_{\Lambda} r=({\cal L}^{\prime} x)_t+\Lambda({\cal L}^{\prime}x)_s =


{\cal L}^{\prime} (x_t+{\cal P} \Lambda{\cal L}^{\prime}x_s) +


D_{\Lambda} {\cal L}^{\prime} \cdot x.


$$


С учетом равенств (1.3), (1.1) и (2.2) вновь получаем систему в инвариантах


$$


r_t+\Lambda(s,t)r_s={\cal L}^{\prime} f({\cal P}r,s,t)+


D_{\Lambda}{\cal L}^{\prime} \cdot {\cal P}r. 


$$


Понятно, что на самом деле эта система отличается от системы (2.3) лишь внешне.


В этом можно легко убедиться с помощью тождества


$$


{\cal L}^{\prime} D {\cal P}+D_{\Lambda} {\cal L}^{\prime} \cdot {\cal P} =


{\cal L}^{\prime} ({\cal P}_t+A{\cal P}_s)+({\cal L}^{\prime}_t +


\Lambda {\cal L}^{\prime}_s){\cal P}=({\cal L}^{\prime} {\cal P})_t+


\Lambda ({\cal L}^{\prime}{\cal P})_s=0.


$$





Обозначим для краткости правую часть инвариантной системы вектор-функцией 


$g(r,s,t)={\cal L}^{\prime} (f({\cal P}r,s,t) - D{\cal P} \cdot r)$.


Очевидно, все аналитические свойства вектор-функций $f$ и $g$ одинаковы.





Дифференциальный оператор $D_{\Lambda}$ инвариантной системы устроен 


существенно проще дифференциального оператора $D$ исходной системы, т.\,к.


ввиду диагональности матрицы $\Lambda$ каждая из $n$ его компонент содержит 


производные по $t$ и по $s$ только от одной неизвестной функции. Примем


обозначения


\begin{equation}


\bigg(\frac{dr_i}{dt}\bigg)_i=r_{i_t}+\lambda_i(s,t)r_{i_s},


\quad i=1,2,\dotsc,n.


\end{equation}


Тогда


$$


D_{\Lambda}r=\bigg(\bigg(\frac{dr_1}{dt} \bigg)_1, \bigg(


\frac{dr_2}{dt}\bigg)_2,\dotsc, \bigg(\frac{dr_n}{dt}\bigg)_n \bigg).


$$








\section{Интегральный эквивалент дифференциальной системы}





Зафиксируем временно произвольную точку $(s,t) \in {\Pi}$ и выпустим из 


нее все характеристики $\xi=s^{(i)}(s,t;\tau)$, $i=1,2,\dotsc,n$, в 


обратном времени $\tau \leq t$. Пусть $(\check{s}^{(i)} (s,t), \check{t}^{(i)}


(s,t))$ --- начальные точки данных характеристик, 


$(\check{s}^{(i)}, \check{t}^{(i)}) \in \overset{<>}{\partial}{}{\Pi} \cup 


\overset{\vee}{\partial}{}{\Pi}$. 


Проинтегрировав уравнения системы (2.3) на отрезках $[\check{t}^{(i)},t]$ 


вдоль соответствующих характеристик $\xi=s^{(i)}(s,t;\tau)$, будем иметь 


\begin{equation}


r_i (s,t)=r_i^o(\check{s}^{(i)},\check{t}^{(i)}) +


\int_{\check{t}^{(i)}}^{t} g_i(r,\xi,\tau) 


\Big|_{\xi=s^{(i)}(s,t; \tau)} d\tau,\quad


i=1,2,\dotsc,n,\quad (s,t) \in {\Pi},


\end{equation}


где $r^o=r^o(s,t)$ --- некоторая определенная на 


$\overset{<>}{\partial}{}{\Pi} \cup \overset{\vee}{\partial}{}{\Pi}$ 


вектор-функция, задающая начально-граничные значения решения $r$.





Для установления эквивалентности дифференциальной (2.3) и интегральной


(3.1) систем на классических (гладких) решениях остается доказать, что решения


интегральной системы (3.1) являются также решениями дифференциальной системы 


(2.3) для любых гладких входных данных $r^o$ и $g$. Доказательство можно 


провести путем непосредственного вычисления производных (2.4) функций 


$r_i (s,t)$, определенных равенствами (3.1). Не приводя здесь этих простых, но 


громоздких выкладок, отметим лишь тождества


\begin{gather}


s_t^{(i)} (s,t;\tau)+ \lambda_i (s,t) s_s^{(i)} (s,t; \tau) \equiv 0,\quad


t_t^{(i)} (s,t;\xi)+\lambda_i (s,t) t_s^{(i)} (s,t; \xi) \equiv 0, \\


%


(s,t) \in {\Pi},\ \ \tau \in T,\ \ \xi \in S,\ \ i=1,2,\dotsc,n,\notag


\end{gather}


благодаря которым обращаются в нуль все слагаемые получающегося в итоге 


выражения за исключением $g_i (r,s,t)$.





Действительно, нетрудно заметить, что в этом выражении левые части тождеств 


(3.2) будут выступать в роли коэффициентов при частных производных по $\xi$


и по $\tau$ соответственно, а единственным слагаемым, не содержащим в 


качестве коэффициента таких частных производных, и будет функция $g_i (r,s,t)$


как результат дифференцирования правой части равенства в системе (3.1) по 


верхнему пределу интегрирования $t$. В свою очередь каждое из тождеств 


(3.2) означает независимость интегральной кривой от выбора ``начальной'' точки


на ней же и может быть получено либо нахождением частных производных решения


обыкновенного дифференциального уравнения по начальным данным [10] в виде


\begin{gather*}


s^{(i)}_t (s,t;\tau)=- \lambda_i (s,t) \exp\bigg(\int_{\tau}^t \lambda_{i_s}


(s^{(i)}(s,t; \alpha), \alpha) d\alpha\bigg), \\


%


s^{(i)}_s (s,t;\tau)=\exp\bigg(\int_{\tau}^t \lambda_{i_s}


(s^{(i)}(s,t; \alpha), \alpha) d\alpha\bigg), \\


%


i=1,2,\dotsc,n;\ \ (s,t) \in {\Pi},\ \ \tau \in T,


\end{gather*}


либо, что проще, дифференцированием тавтологий


\begin{gather*}


\xi \equiv s^{(i)}(s^{(i)}(\xi, \tau;t),t;\tau),\quad


\tau \equiv t^{(i)}(s,t^{(i)}(\xi, \tau;s);\xi),\\


%


(\xi,\tau) \in {\Pi},\ \ t \in T,\ \ s \in S,\ \ i=1,2,\dotsc,n,


\end{gather*}


по переменным $t$ и $s$ соответственно.





Итак, дифференциальная (2.3) и интегральная (3.1) системы эквивалентны 


в случае, если они имеют классическое, т.\,е. непрерывное и непрерывно


дифференцируемое в $\Pi$ решение $r=r(s,t)$. Как известно (напр.,


[1], [2]), для этого необходимо требовать аналогичную гладкость от правой части 


системы (1.1). В то же время интегральная система (3.1), очевидно, имеет смысл


не только для гладких вектор-функций $f$, что и позволяет определять 


соответствующие понятия обобщенного решения инвариантной системы (2.3), а 


следовательно, и исходной гиперболической системы (1.1).








\section{Построение корректных граничных условий}





Вначале построим корректные начально-граничные условия для решения инвариантной 


системы (2.3), а затем трансформируем их с помощью тождеств (2.2) в условия 


для решения $x$ исходной гиперболической системы (1.1).





Интегральный эквивалент (3.1) инвариантной системы (2.3) использует граничные 


значения инвариантов Римана $r_i$ только в начальных точках 


характеристик. Следовательно, граничные условия должны быть устроены так, 


чтобы они ``закрепляли'' каждый инвариант Римана в начальных точках и,


напротив, не фиксировали его в конечных точках  ``своих''  характеристик. 


Несмотря на однозначность этого положения, оно допускает два различных способа


формальной записи граничных условий. Один из них


(напр., [1], [2], [5], [7] ,[11]) 


жестко привязан к понятиям  ``левая'', ``нижняя'', ``правая''  границы области 


$\Pi$ и требует строгого знакопостоянства собственных значений 


$\lambda_i$, $i=1,2,\dotsc,n$, всюду в области $\Pi$. В этом случае вектор 


инвариантов Римана разбивается на три подвектора по признаку 


$\lambda_i < 0$, $\lambda_i =0$ или $\lambda_i > 0$. Инварианты


Римана $r_i$, для которых $\lambda_i > 0$, закрепляются на левой границе


$\overset{<}{\partial}{}{\Pi}$, инварианты Римана $r_i$, соответствующие 


$\lambda_i < 0$, --- на правой границе $\overset{>}{\partial}{}{\Pi}$, 


наконец, на нижней границе $\overset{\vee}{\partial}{}{\Pi}$ закрепляются все


инварианты $r_i$ вне зависимости от знака $\lambda_i$.





Другой способ записи граничных условий (напр., [12]--[17]) не требует 


указанных ограничений, более компактен, имеет одинаковую форму записи как в 


одномерных, так и в многомерных гиперболических системах и применим для 


произвольных односвязных областей с кусочно-гладкими границами.





Введем вспомогательные функции, которые будут использоваться не 


только здесь при постановке начально-граничных условий, но и далее для других


целей. Поэтому определение данных функций дадим для произвольной односвязной


области $Q$ с кусочно-гладкой границей $\partial Q$. Пусть 


$\wh \nu=(\nu_o,\nu)$ --- единичный вектор внешней нормали к границе 


$\partial Q$ в точках $(s,t) \in \partial Q$, причем $\nu_o=\nu_o(s,t)$ --- 


величина проекции вектора $\wh \nu (s,t)$ на ось $t$,


а $\nu=\nu (s,t)$ на ось $s$.





Определим функции $z_i=z_i (\partial Q;s,t)$ равенствами


\begin{equation}


z_i(\partial Q;s,t)=\nu_o (s,t)+\lambda_i (s,t) \nu (s,t),


\quad i=1,2,\dotsc,n;\quad (s,t) \in \partial Q.


\end{equation}





Очевидно, значение функции $z_i$ в точке $(s,t)$ есть значение скалярного 


произведения вектора внешней нормали $\wh \nu (s,t)$ и вектора 


$(1, \lambda_i (s,t))$, который в силу равенства 


$(dt,ds)=(1, \lambda_i (s,t)) dt$ лежит на касательной, проведенной в точке


$(s,t)$ к характеристике $s^{(i)} (s,t; \cdot)$. Отсюда справедлива следующая 


классификация точек $(s,t) \in \partial Q$ относительно характеристики


$s^{(i)} (s,t; \cdot)$. Если $z_i < 0$, то $(s,t) \in \partial Q$


является начальной точкой этой характеристики, т.\,к. для любого достаточно 


малого $\varepsilon > 0$ точки $(s^{(i)} (s,t; \tau), \tau) \in \Int Q$


при всех $\tau \in (t,t+\varepsilon)$. Если $z_i > 0$, то 


$(s,t) \in \partial Q$ --- конечная точка характеристики $s^{(i)} (s,t; \cdot)$,


т.\,к. для любого достаточно малого $\varepsilon > 0$ точки


$(s^{(i)} (s,t; \tau), \tau) \in \Int Q$ при всех 


$\tau \in (t-\varepsilon, t)$. Случай $z_i=0$ означает, что в точке


$(s,t)$ реализуется либо прикосновение к границе $\partial Q$, либо скольжение 


вдоль нее характеристики $s^{(i)} (s,t; \cdot)$. При этом точка $(s,t)$ 


может оказаться начальной, конечной или внутренней точкой характеристики


$s^{(i)} (s,t; \cdot)$. Однако ввиду экзотичности этих точек (в том смысле, 


что на границе $\partial Q$ мера множества всех начальных и конечных точек,


соответствующих равенству $z_i=0$, равна нулю) их детальная классификация 


не потребуется, хотя и достаточно очевидна.





С помощью функций (4.1) определим функции со значением 0 или 1


$$


z_i^{+}=\sign (z_i) (1+\sign (z_i))/2,\quad


z_i^{-}=\sign (z_i) (\sign (z_i) - 1)/2,\quad i=1,2,\dotsc,n,


$$


и матричные функции 


$$


Z=\diag \{ z_1, z_2,\dotsc, z_n \},\quad


Z^{+}= \diag \{ z_1^{+}, z_2^{+},\dotsc, z_n^{+} \},\quad


Z^{-}= \diag \{ z_1^{-}, z_2^{-},\dotsc, z_n^{-} \}.


$$





Пусть теперь 


$z_i=z_i (\partial {\Pi}; s,t)$, $i=1,2,\dotsc,n$; 


$(s,t) \in \partial {\Pi}$, а $r^o=r^o (s,t)$ --- $n$-мерная вектор-функция, 


заданная на границе $\partial {\Pi}$. Тогда начально-граничные условия 


для инвариантной системы (2.3) можно записать в векторно-матричном виде


\begin{equation}


Z^{-} (s,t) (r(s,t) - r^o (s,t))=0,\quad (s,t) \in \partial {\Pi}.


\end{equation}


Эти условия пригодны не для всех приложений 


гиперболических систем


уравнений. Дело в том, что на практике, как правило, граничные 


условия формулируются не для линейных комбинаций ${\cal L}^{\prime}x$ искомого


решения $x$, т.\,е. фактически инвариантов Римана $r$, а непосредственно для 


компонент вектор-функции $x$, которые только лишь и имеют конкретный 


содержательный смысл. Их расшифровка в терминах инвариантов Римана приводит к 


более общим, чем (4.2), так называемым cмешанным начально-краевым 


условиям.





Пусть $q=q (x,s,t)$ --- заданная $n$-мерная вектор-функция с областью 


определения $R^n \times \partial \Pi$. Тогда вместо начально-граничных 


условий (4.2) можно поставить смешанные условия 


\begin{equation}


Z^{-} (s,t) [ r(s,t) - q (Z^{+} (s,t) r(s,t),s,t) ]=0.


\end{equation}


Условия (4.3) обобщают условия (4.2) и совпадают с ними,


лишь если $q(Z^{+} r,s,t)=r^o (s,t)$. 


Отметим, что вопросы обоснования корректности смешанных задач гиперболического 


типа изучались многими авторами (напр., [1], [2], [5], [7]), но лишь для систем, 


записанных в инвариантах Римана и с линейными по решению 


(за исключением [7]) функциями $q$.





После замены в соотношении (4.3) инвариантов Римана $r$ на исходную функцию 


$x$ по правилу (2.2) будем иметь 


\begin{equation}


Z^{-} (s,t) [ {\cal L}^{\prime}(s,t)x(s,t) - q (Z^{+} (s,t) {\cal L}^{\prime}


(s,t) x(s,t),s,t) ]=0.


\end{equation}








\section{Обобщенное решение}





В основу определения обобщенного решения положим идею, которая для этой цели 


часто используется как в системах обыкновенных дифференциальных уравнений


[18], [19], так и в гиперболичсеких системах


[1], [11], [16], [17]. Именно,


обобщенным решением системы дифференциальных уравнений объявляется 


решение эквивалентной ей системы интегральных уравнений. Понятно, что 


эквивалентность систем устанавливается на классических (гладких) решениях. Как


было уже показано, интегральная система (3.1) эквивалентна дифференциальной 


начально-краевой задаче в инвариантах (2.3), (4.2). В точности так же можно 


доказать эквивалентность смешанной задачи (2.3), (4.3) и интегральной системы


\begin{equation} 


r_i (s,t)=q_i (Z^{+}r, \check{s}^{(i)}, \check{t}^{(i)}) +


\int_{\check{t}^{(i)}}^{t}


g_{i} (r, \xi, \tau) \Big|_{\xi=s^{(i)}(s,t; \tau)}


d\tau,\quad i=1,2,\dotsc, n,\quad (s,t) \in {\Pi}.


\end{equation}


Поэтому интегральный эквивалент смешанной задачи (1.1), (4.4) имеет вид


\begin{equation} 


x (s,t)=\sum_{i=1}^n p^{(i)} (s,t) \bigg[ q_i (Z^{+}


{\cal L}^{\prime} x, \check{s}^{(i)}, \check{t}^{(i)}) +


\int_{\check{t}^{(i)}}^{t}g_{i} ({\cal L}^{\prime}x, \xi, \tau)


\Big|_{\xi=s^{(i)}(s,t; \tau)} d\tau \bigg], 


\ \ i=1,2,\dotsc, n,\ \ (s,t) \in {\Pi},


\end{equation}


что позволяет дать





\begin{defns}


Обобщенным решением гиперболической системы уравнений 


(1.1) с на\-чаль\-но-гра\-нич\-ны\-ми условиями (4.4) назовем вектор-функцию 


$x=x(s,t)$, которая почти всюду в $\Pi$ удовлетворяет равенству (5.2).


\end{defns}





Прежде чем перейти непосредственно к формулировкам и доказательствам 


соответствующих утверждений, введем вспомогательные конструкции.





Пусть $G(s,t) \subset {\Pi}$ --- область определения решения 


$x$, соответствующая точке $(s,t) \in {\Pi}$. Левая 


$\xi=s_{\min} (s,t; \tau)$ и правая $\xi=s_{\max} (s,t; \tau)$ 


границы области $G(s,t)$ находятся как решения обыкновенных дифференциальных 


уравнений


$$


\frac{d\xi}{d\tau}=


\begin{cases}


\max\limits_{i=1,2,\dotsc,n} \lambda_i(\xi,\tau), 


& (\xi,\tau) \in (s_o,s_1]; \\


0, & \xi=s_o, 


\end{cases}


\qquad


\frac{d\xi}{d\tau}=


\begin{cases}


\min\limits_{i=1,2,\dotsc,n} \lambda_i(\xi,\tau), 


& (\xi,\tau) \in [s_o,s_1); \\


0, & \xi=s_1, 


\end{cases}


$$


соответственно, которые интегрируются в ``обратном времени''


$\tau < t$ на отрезке $T_G (s,t)=[\check{t}(s,t),t]$, где


момент $\check{t}(s,t)$ окончания интегрирования либо служит корнем уравнения 


$s_{\min} (s,t; \tau)=s_{\max} (s,t; \tau)$, либо совпадает с $t_0$.


Поэтому облаcть $G (s,t)$ определяется по правилу


$$


G(s,t)=\{ (\xi, \tau) \in {\Pi}: 


s_{\min} (s,t; \tau) < \xi < s_{\max} (s,t; \tau),


\ \tau \in T_G (s,t) \}.


$$





Укажем теперь специальные функциональные пространства, которым (это будет 


показано далее) при соответствующем выборе вектор-функций $f$ и $q$ 


принадлежат решения $r$ и $x$ интегральных систем (5.1) и (5.2). Из 


предыдущего параграфа следует, что для почти каждой точки $(s,t) \in {\Pi}$ 


и каждого $i=1,2,\dotsc,n$ существует единственная пара точек 


$(\check{s}^{(i)}, \check{t}^{(i)})$, $(\wh{s}^{(i)}, \wh{t}^{(i)}) \in 


\partial {\Pi}$. При этом


$z_i (\partial {\Pi}; \check{s}^{(i)}, \check{t}^{(i)}) < 0$,


$z_i (\partial {\Pi}; \wh{s}^{(i)}, \wh{t}^{(i)}) > 0$, т.\,е.


$(\check{s}^{(i)}, \check{t}^{(i)})$ --- начальная, 


$(\wh{s}^{(i)}, \wh{t}^{(i)})$ --- конечная точки характеристики 


$\xi=s^{(i)}(s,t; \tau)$, 


$\tau \in [\check{t}^{(i)}, \wh{t}^{(i)}]=T_i(s,t)$. 


Положим


\begin{align*}


\Lambda_{p}^n ({\Pi})&= \{ r \in {\mathrm L}_{p}^n


({\Pi}): r_i (s^{(i)} 


(s,t; \cdot), \cdot) \in {\mathrm{AC}}(T_i (s,t)),


\ i=1,2,\dotsc,n, \ (s,t) \in {\Pi} \}, \\


%


{\mathrm A}_{p}^n ({\Pi})&= \{ x \in {\mathrm L}_{p}^n ({\Pi}): 


x(s,t)={\cal P}(s,t) r(s,t),\ (s,t) \in {\Pi}, 


\ r \in \Lambda_{p}^n ({\Pi}) \}.


\end{align*}


Здесь ${\mathrm L}_{p}^n ({\Pi})$ ---


пространство $n$-мерных вектор-функций, 


суммируемых со степенью $p$, $p \geq 1$, в области $\Pi$


по Лебегу, ${\mathrm {AC}}(T)$ -- пространство абсолютно непрерывных на отрезке


$T$ функций. Очевидно, что $x \in {\mathrm A}_{p}^n ({\Pi})$ тогда и только


тогда, когда ${\cal L}^{\prime}x \in \Lambda_{p}^n ({\Pi})$.





\begin{thms}


Пусть вектор-функции $f=f(x,s,t)$, $q=q(x,s,t)$ в 


областях $R^n \times {\Pi}$ и $R^n \times \partial {\Pi}$ соответственно 


удовлетворяют условию Липшица по переменной $x$ при фиксированных $(s,t)$ и 


$f(x,\cdot,\cdot) \in {\mathrm L}_{p}^n ({\Pi})$, 


$q(x,\cdot,\cdot) \in {\mathrm L}_{p}^n (\partial {\Pi})$ 


при фиксированных $x$. 


Тогда существует решение $x \in {\mathrm A}_{\mathrm P}^n ({\Pi})$ системы 


интегральных уравнений $(5.2)$.


\end{thms}





{\bf Доказательство} теоремы проводится по следующей схеме. Строится 


последовательность функций $\{ x^{(k)} \}$ методом последовательных 


приближений


$$ 


x^{(k+1)} (s,t)=\sum_{i=1}^n p^{(i)} (s,t) \bigg[ q_i (Z^{+}


{\cal L}^{\prime} x^{(k+1)}, \check{s}^{(i)}, \check{t}^{(i)}) +


\int_{\check{t}^{(i)}}^{t} 


g_{i} ({\cal L}^{\prime}x^{(k)}, \xi, \tau) \Big|_{\xi=


s^{(i)}(s,t; \tau)} d\tau \bigg]. 


$$


Показывается, что отображение, генерирующее последовательность 


$\{ x^{(k)} \}$, является сжимающим. Следовательно, существует предельная 


точка $x \in {\mathrm A}_{p}^{n} ({\Pi})$, удовлетворяющая интегральной 


системе уравнений (5.2).





На основе свойства сжимаемости отображения, фигурирующего в правой части


системы (5.2), доказывается 





\begin{thms}


Пусть $\wt{x}$ и $x$ --- обобщенные решения смешанных


задач $(1.1)$, $(4.4)$, соответствующие входным данным $\wt{f}$, $\wt{q}$ и


$f$, $q$, которые удовлетворяют условиям теоремы $5.1$. Тогда справедлива оценка


\begin{gather*}


\| \wt{x} (s,t) - x(s,t)\| \leq {K} \bigg[


\sum\limits_{i=1}^{n} \bigg(\| \triangle q (Z^{+}{\cal L}^{\prime}x, 


\check{s}^{(i)},\check{t}^{(i)}) \|+


\int\limits_{\check{t}^{(i)}}^{t} \|\triangle f ({\cal L}^{\prime}x, \xi, \tau)


\| \Big|_{\xi=s^{(i)}(s,t; \tau)} d\tau \bigg) + \\


%


+\int\limits_{\partial G(s,t) \cap \partial {\Pi}} 


\| \triangle q (Z^{+} {\cal L}^{\prime}x, \xi,\tau) \| d\omega


+\iint\limits_{G(s,t)}


\| \triangle f ({\cal L}^{\prime}x, \xi,\tau) \| d\xi\, d\tau \bigg],


\end{gather*}


в которой $\triangle q (x,s,t)=\wt{q} (x,s,t) - q(x,s,t)$,


$\triangle f (x,s,t)=\wt{f} (x,s,t) - f(x,s,t)$ --- разности входных 


данных$;$ $d\omega=\sqrt{d\xi^2+d\tau^2}$ --- дифференциал длины границы


$\partial {\Pi}$ в интеграле первого рода$;$ $K$ --- положительная 


константа, зависящая только от констант Липшица вектор-функций 


$\wt{f}$, $\wt{q}$, $f$, $q$ и размера отрезка $T$.


\end{thms}





В силу теоремы 5.2 обобщенное решение смешанной задачи (1.1),


(4.4) является единственным и непрерывно зависит от входных данных.





В заключение обоснуем справедливость формулы интегрирования по частям. Для 


этого введем функциональные пространства, сопряженные


с пространствами $\Lambda_{p}^n ({\Pi})$ и $\mathrm A_{p}^n ({\Pi})$


соответственно,


\begin{align*} 


\Lambda^{*n}_m ({\Pi})&=\big\{ \varphi \in \mathrm L^n_m ({\Pi}):


\varphi_i (s^{(i)} (s,t; \cdot), \cdot) \in {\mathrm{AC}} (T_i (s,t)),


\ i=1,2,\dotsc,n;\ (s,t) \in {\Pi} \big\}, \\


%


{\mathrm A}^{*n}_m ({\Pi})&=\big\{ \psi \in \mathrm L^n_m ({\Pi}):


\psi (s,t)={\cal L} (s,t) \varphi (s,t),\ (s,t) \in {\Pi},\ \varphi \in 


\Lambda^*{}^n_m ({\Pi}) \big\},


\end{align*}


где $m,p \geq 1$ и ${1}/{m}+{1}/p=1$.





Заметим, что $\psi (s,t) \in {\mathrm A}^{*n}_m ({\Pi})$, если и только


если ${\cal P}^{\prime}\psi \in \Lambda^{*n}_m ({\Pi})$.





Ввиду абсолютной непрерывности функций $\varphi_i$ вдоль характеристик


``своих'' семейств имеет смысл почти всюду в $\Pi$ дифференциальный оператор


${\mathrm D}^*_{\Lambda}$, определенный на функциях 


$\varphi \in \Lambda^{*n}_m (\Pi)$ по правилу


$$


{\mathrm D}^*_{\Lambda} \varphi={\mathrm D}_{\Lambda} \varphi+\Lambda_s (s,t) 


\varphi.


$$





По схеме работы [5] можно доказать, что для произвольной односвязной 


подобласти $Q \subseteq {\Pi}$ с кусочно-гладкой границей $\partial Q$ и


любых вектор-функций $r \in \Lambda^n_{p} (\Pi)$,


$\varphi \in \Lambda^{*n}_m (\Pi)$ справедлива формула интегрирования по 


частям


\begin{equation}


\iint\limits_{Q}


\langle \varphi, {\mathrm D}_{\Lambda} r \ra ds\, dt=\int\limits_{\partial Q} 


\la \varphi, Z(s,t) r \ra d \omega - \iint\limits_{Q} 


\la {\mathrm D}^*_{\Lambda} \varphi, r\ra ds\, dt,


\end{equation}


в которой $\la\ ,\ \ra$ --- знак скалярного произведения в пространстве $R^n$;


$Z =\diag \{ z_1, z_2,\dotsc, z_n \}$, а функции $z_i=z_i (s,t)$, 


$i=1,2,\dotsc,n$, определены по правилу (4.1).





В отличие от вектор-функций пространств $\Lambda^n_{p} ({\Pi})$ и 


$\Lambda^{*n}_m ({\Pi})$, $i$-е координаты которых являются абсолютно 


непрерывными функциями вдоль характеристик $i$-го семейства, вектор-функции


пространств ${\mathrm A}^n_{p} ({\Pi})$ и ${\mathrm A}^{*n}_m ({\Pi})$ 


устроены более сложно. В частности (и это можно проиллюстрировать конкретными


примерами) компоненты $x_i (\psi_i)$ вектор-функций $x (\psi)$ могут


не только не иметь производных, но и быть разрывными в каждой точке 


$(s,t) \in {\Pi}$ по любому направлению. Тем не менее, корректный


способ вычисления значения оператора $\mathrm D$ в точках 


$x \in \Lambda^n_{p} ({\Pi})$ нетрудно выписать с помощью оператора


$\mathrm D_{\Lambda}$. Достаточно положить 


\begin{equation}


{\mathrm D}x={\mathrm D}{\cal P} \cdot {\cal L}^{\prime} x+


{\cal P}{\mathrm D}_{\Lambda} ({\cal L}^{\prime} x).


\end{equation}


Аналогично поступим и с сопряженным оператором ${\mathrm D}^*$, который на 


гладких функциях $\psi$ задается равенством


${\mathrm D} \psi^{*}=\psi_{t}+({\mathrm A}^{\prime} \psi)_s$. В точках 


$\psi \in {\mathrm A}^{*n}_m({\Pi})$ определим его по правилу


\begin{equation}


{\mathrm D}^* \psi={\mathrm D}^* {\cal L} \cdot {\cal P}^{\prime} \psi - 


{\cal L} \Lambda_s \psi+{\cal L} {\mathrm D}^*_{\Lambda} 


({\cal P}^{\prime} \psi).


\end{equation}





\begin{thms}


Пусть $x \in {\mathrm A}^n_p ({\Pi})$,


$\psi \in {\mathrm A}^{*n}_m ({\Pi})$. Тогда для любой односвязной подобласти


$Q \subseteq \Pi$ с кусочно-гладкой границей $\partial Q$ справедлива


формула интегрирования по частям


$$


\iint\limits_{Q} 


\la\psi, Dx\ra ds\, dt=\int\limits_{\partial Q}


\la\psi, {\mathrm B} x\ra d \omega -


\int\limits_{Q} \la D^* \psi, x\ra ds\, dt,


$$ 


где матричная функция ${\mathrm B=\mathrm B}(s,t)$ вычисляется на границе 


$\partial Q$ по формуле


$$


{\mathrm B}(s,t)=\nu_o (s,t) {\mathrm E}+\nu (s,t) {\mathrm A}(s,t),


$$


в которой $\wh\nu=(\nu_o, \nu)$ -- единичная внешняя нормаль к границе 


$\partial Q$.


\end{thms}





Доказательство теоремы основано на равенстве (5.3) и формулах (5.4), (5.5).








\begin{thebibliography}{99}





\bibitem{1}


Рождественский Б.Л., Яненко Н.Н. {\em Системы квазилинейных уравнений


и их приложения к газовой динамике}. -- М.: Наука, 1978. -- 687~с.


\bibitem{2}


Годунов С.К. {\em Уравнения математической физики}. -- М.: Наука, 1979. -- 


392~с.


\bibitem{3}


Егоров А.И., Знаменская Л.Н. {\em Управление упругими колебаниями $($обзор$)$} //


Оптимизация, управление, интеллект. -- Иркутск: ИД СТУ СО РАН, 2000. -- \No\,5.


-- С.\,249--257.


\bibitem{4}


Москаленко А.И., Москаленко Р.А., Хамитов Г.П. {\em Моделирование возрастной 


структуры производственных фондов с точки зрения задач управления 


экономическими системами} // Оптимизация, управление, интеллект. -- Иркутск: 


ИД СТУ СО РАН, 2000. -- \No\,5. -- С.\,418--427.


\bibitem{5}


Потапов М.М. {\em Обобщенное решение смешанной задачи для полулинейной


гиперболической системы первого порядка} // Дифференц. уравнения. -- 1983.--


Т.\,19. -- \No\,10. -- С.\,1826--1828.


\bibitem{6}


Морозов С.Ф. {\em Начально-краевая задача для почти линейной гиперболической 


системы на плоскости} // Краевые задачи. -- Пермь, 1989. -- С.\,141--149.


\bibitem{7}


Матвеев А.С. {\em Обобщенные решения полулинейной системы уравнений в частных


производных гиперболического типа и задачи управления} // Деп. в ВИНИТИ 


23.07.90, \No\,2983-80. -- 39~с.


\bibitem{8}


Ланкастер П. {\em Теория матриц}. -- М.: Наука, 1982. -- 270 с.


\bibitem{9}


Levy H. {\em \"Uber Anfangswertproblem f\"ur eine hyperbolische 


nichtlineare Differentialgleichungen zweiter Ordnung mit zwei unabh\"angigen


Ver\"anderlichen} // Math. Ann. -- 1927. -- V.\,97. -- P.\,179--191.


\bibitem{10}


Петровский И.Г. {\em Лекции по теории обыкновенных дифференциальных 


уравнений}. -- М.: Наука, 1964. -- 272~с.


\bibitem{11}


Васильев О.В., Срочко В.А., Терлецкий В.А. {\em Методы оптимизации и


их приложения}. Ч.\,2. {\em Оптимальное управление}. -- Новосибирск: Наука,


1990. -- 151~с.


\bibitem{12}


Friedrichs K.O. {\em Symmetric positive linear differential equations}


// Comm. Pure Appl. Math. -- 1958. -- V.\,11. -- \No\,3. -- P.\,333--418.


\bibitem{13}


Lax P.D., Phillips R.S. {\em Local boundary conditions for dissipative 


symmetric linear differential operators} // Comm. Pure Appl. Math. -- 1960.


-- V.\,13. -- \No\,3. -- P.\,427--455.


\bibitem{14}


Phillips R.S., Sarason L. {\em Singular symmetric first order


differential operators} // J. Math. and Mech. -- 1966. -- V.\,15.


-- \No\,2. -- P.\,235--271.


\bibitem{15}


Vasiliev O.V., Terletsky V.A., Arguchintsev A.V. {\em Iterative


processes in optimization of semilinear hyperbolic system} //


11-th IFAC World Congress, V.\,6. -- Tallinn, 1990. -- P.\,216--220.


\bibitem{16}


Терлецкий В.А. {\em Вариационный принцип максимума в управляемых системах


одномерных гиперболических уравнений} // Изв. вузов. Математика.


-- 1999. -- \No\,12. -- С.\,82--90.


\bibitem{17}


Терлецкий В.А. {\em Обобщенное решение многомерных полулинейных 


гиперболических систем} // Изв. вузов. Математика.


-- 2001. -- \No\,\,12. -- С.\,68--76.


\bibitem{18}


Сансоне Дж. {\em Обыкновенные дифференциальные уравнения}. T.\,II.


-- М.: Ин. лит., 1954. -- 415~с.


\bibitem{19}


Коддингтон Э.А., Левинсон Н. {\em Теория обыкновенных дифференциальных


уравнений}. -- М.: Ин. лит., 1958. -- 474~с.





\end{thebibliography}





\vskip 2cm





\post{\it Иркутский государственный}{\it Поступила}


\post{\it университет}{27.09.2004}





\label{end}


\end{document}


